\scanpage{275}
\uline{Définition} $X$ préschéma sur $S$. Un élément $\xi$ de $\operatorname{Pic}(X/S)$ est dit \uline{numériquement équivalent à $0$} si sa restriction à tte fibre $X_s$ est numériquement \unsure{équivalent} à zéro [ie. un multiple d'icelle peut se déformer dans $0$]. Deux \unsure{éléme} $\xi,\xi'$ de $\operatorname{Pic}(X/S)$ sont dits numériquement équivalents si $\xi-\xi'$ est numériquement équivalent à $0$,

Soit $\operatorname{Pic}^{\tau}(X/S)\subset\operatorname{Pic}^{\deleted{\unsure{\tau}}}(X/S)$ des $\xi$ numériquement \unsure{équivalts} à $0$, et $\deleted{\unsure{\operatorname{Pol}}}\inserted{\unsure{\operatorname{N\acute{e}r}}}'(X/S)=\operatorname{Pic}(X/S)/\operatorname{Pic}^{\tau}(X/S)$. Noter : si \deleted{$S'\to S$ est fid plat [et \unsure{qu. cpt}] et} $S'\to S$ est un \unsure{changt} de base, alors $\operatorname{Pic}(X/S)\to\operatorname{Pic}(X'/S')$ induit $\operatorname{Pic}^{\tau}(X/S)\to\operatorname{Pic}^{\tau}(X'/S')$, et si $S'\to S$ est surjectif, alors \deleted{un élément $\xi$ de} $\operatorname{Pic}^{\tau}(X/S)$ est exactement l'image inverse de $\operatorname{Pic}^{\tau}(X'/S')$ par $\operatorname{Pic}^{\deleted{\unsure{\tau}}}(X/S)\to\operatorname{Pic}(X'/S')$, donc dans ce cas
\[
\deleted{\unsure{\operatorname{Pol}}}\inserted{\unsure{\operatorname{N\acute{e}r}}}'(X/S)\longrightarrow\deleted{\unsure{\operatorname{Pol}}}\inserted{\unsure{\operatorname{N\acute{e}r}}}'(X'/S')\quad\text{est \uline{injectif}}
\]
\margin{\diagram{\begin{tikzcd}X\arrow[d]&X'\arrow[l]\arrow[d]\\S&S'\arrow[l]\end{tikzcd}}}

Considérons $S'\mapsto\operatorname{Pol}'(X'/S')$ comme foncteur contravariant $(\operatorname{Sch})^0_{/S}\longrightarrow(\operatorname{Ens})$, et \uline{rendons le exact pour les morphismes fid plats \deleted{\unsure{qu. cpts}} \inserted{loc. de type} fini} [donc pour la localisation habituelle aussi]. On trouve un foncteur $S'\mapsto\deleted{\unsure{\operatorname{Pol}}}\inserted{\unsure{\operatorname{N\acute{e}r}}}(X'/S')$, et en particulier on a un \inserted{ens} $\deleted{\unsure{\operatorname{Pol}}}\inserted{\unsure{\operatorname{N\acute{e}r}}}(X/S)$, ens des polarisations de $X$ sur $S$.
\mspage{274}
